\documentclass[UTF8, 12pt, fontset=fandol, oneside]{ctexart}
\usepackage{researchnotes}

\title{强化学习：导论}
\author{}
\date{}

\begin{document}

\maketitle

\section*{符号汇总}

大写字母用于随机变量，小写字母用于随机变量的取值以及标量函数。要求为实值向量的量用小写粗体表示（即使它们是随机变量）。矩阵用大写粗体表示。

\begin{longtable}{p{0.22\textwidth}p{0.68\textwidth}}
$\doteq$ & 按定义成立的等式关系 \\
$\approx$ & 近似相等 \\
$\propto$ & 正比于 \\
$\Pr\{X=x\}$ & 随机变量 $X$ 取值为 $x$ 的概率 \\
$X\sim p$ & 随机变量 $X$ 从分布 $p(x)\doteq\Pr\{X=x\}$ 中抽取 \\
$\mathbb{E}[X]$ & 随机变量 $X$ 的期望，即 $\mathbb{E}[X]\doteq\sum_x p(x)x$ \\
$\arg\max_a f(a)$ & 使 $f(a)$ 取得最大值的 $a$ 的一个取值 \\
$\ln x$ & $x$ 的自然对数 \\
$e^x$ & 自然对数的底数 $e\approx2.71828$ 的 $x$ 次幂；$e^{\ln x}=x$ \\
$\mathbb{R}$ & 实数集 \\
$f:\mathcal{X}\to\mathcal{Y}$ & 从集合 $\mathcal{X}$ 的元素到集合 $\mathcal{Y}$ 的元素的函数 $f$ \\
$\leftarrow$ & 赋值 \\
$(a,b]$ & $a$ 与 $b$ 之间的实区间，包含 $b$ 但不包含 $a$ \\
\end{longtable}

\noindent\textbf{在多臂赌博机问题中：}
\begin{longtable}{p{0.22\textwidth}p{0.68\textwidth}}
$\varepsilon$ & 在 $\varepsilon$-贪心策略中采取随机动作的概率 \\
$\alpha,\beta$ & 步长参数 \\
$\gamma$ & 折扣率参数 \\
$\lambda$ & 资格迹的衰减率参数 \\
$\mathbf{1}_{predicate}$ & 指示函数（若 $predicate$ 为真，则 $\mathbf{1}_{predicate}\doteq1$，否则为 $0$） \\
$k$ & 动作（臂）的个数 \\
$t$ & 离散时间步或游玩次数 \\
$q_*(a)$ & 动作 $a$ 的真实价值（期望奖励） \\
$Q_t(a)$ & 在时刻 $t$ 对 $q_*(a)$ 的估计 \\
$N_t(a)$ & 在时刻 $t$ 之前动作 $a$ 被选择的次数 \\
$H_t(a)$ & 在时刻 $t$ 学到的、选择动作 $a$ 的偏好 \\
$\pi_t(a)$ & 在时刻 $t$ 选择动作 $a$ 的概率 \\
$\bar R_t$ & 给定 $\pi_t$ 时，对期望奖励在时刻 $t$ 的估计 \\
\end{longtable}

\noindent\textbf{在马尔可夫决策过程中：}
\begin{longtable}{p{0.22\textwidth}p{0.68\textwidth}}
$s,s'$ & 状态 \\
$a$ & 一个动作 \\
$r$ & 一个奖励 \\
$\mathcal{S}$ & 所有非终止状态的集合 \\
$\mathcal{S}^+$ & 所有状态的集合，包括终止状态 \\
$\mathcal{A}(s)$ & 状态 $s$ 下所有可用动作的集合 \\
$\mathcal{R}$ & 所有可能奖励的集合，是 $\mathbb{R}$ 的一个有限子集 \\
$\subset$ & 子集；例如 $\mathcal{R}\subset\mathbb{R}$ \\
$\in$ & 属于；例如 $s\in\mathcal{S}$、$r\in\mathcal{R}$ \\
$|\mathcal{S}|$ & 集合 $\mathcal{S}$ 中元素的个数 \\
$t$ & 离散时间步 \\
$T,T(t)$ & 一个幕的最后时间步，或包含时间步 $t$ 的那个幕的最后时间步 \\
$A_t$ & 时刻 $t$ 的动作 \\
$S_t$ & 时刻 $t$ 的状态，通常由 $S_{t-1}$ 和 $A_{t-1}$ 随机地产生 \\
$R_t$ & 时刻 $t$ 的奖励，通常由 $S_{t-1}$ 和 $A_{t-1}$ 随机地产生 \\
$\pi$ & 策略（决策规则） \\
$\pi(s)$ & 在确定性策略 $\pi$ 下，于状态 $s$ 采取的动作 \\
$\pi(a|s)$ & 在随机性策略 $\pi$ 下，于状态 $s$ 采取动作 $a$ 的概率 \\
$G_t$ & 时刻 $t$ 之后的回报 \\
$h$ & 视界，即在前向视图中向前看到的时间步 \\
$G_{t:t+n},G_{t:h}$ & 从 $t+1$ 到 $t+n$ 或到 $h$ 的 $n$ 步回报（经过折扣与校正） \\
$\bar G_{t:h}$ & 从 $t+1$ 到 $h$ 的平坦回报（未折扣且未校正）（第 5.8 节） \\
$G_t^\lambda$ & $\lambda$-回报（第 12.1 节） \\
$G_{t:h}^\lambda$ & 截断并校正的 $\lambda$-回报（第 12.3 节） \\
$G_t^{\lambda s},G_t^{\lambda a}$ & 分别由估计的状态价值或动作价值校正的 $\lambda$-回报（第 12.8 节） \\
$p(s',r|s,a)$ & 从状态 $s$ 采取动作 $a$ 后，转移到状态 $s'$ 并获得奖励 $r$ 的概率 \\
$p(s'|s,a)$ & 从状态 $s$ 采取动作 $a$ 后，转移到状态 $s'$ 的概率 \\
$r(s,a)$ & 在状态 $s$ 采取动作 $a$ 后的期望即时奖励 \\
$r(s,a,s')$ & 在动作 $a$ 下从 $s$ 转移到 $s'$ 时的期望即时奖励 \\
$v_\pi(s)$ & 策略 $\pi$ 下状态 $s$ 的价值（期望回报） \\
$v_*(s)$ & 最优策略下状态 $s$ 的价值 \\
$q_\pi(s,a)$ & 策略 $\pi$ 下在状态 $s$ 采取动作 $a$ 的价值 \\
$q_*(s,a)$ & 最优策略下在状态 $s$ 采取动作 $a$ 的价值 \\
$V,V_t$ & 对状态价值函数 $v_\pi$ 或 $v_*$ 的数组估计 \\
$Q,Q_t$ & 对动作价值函数 $q_\pi$ 或 $q_*$ 的数组估计 \\
$\bar V_t(s)$ & 期望近似动作价值，例如 $\bar V_t(s)\doteq\sum_a\pi(a|s)Q_t(s,a)$ \\
$U_t$ & 时刻 $t$ 用于估计的目标 \\
$\delta_t$ & 时刻 $t$ 的时序差分（TD）误差（一个随机变量）（第 6.1 节） \\
$\delta_t^s,\delta_t^a$ & TD 误差的状态特定形式和动作特定形式（第 12.9 节） \\
$n$ & 在 $n$ 步方法中，$n$ 是自举的步数 \\
$d$ & 维数，即 $\mathbf{w}$ 的分量个数 \\
$d'$ & 另一维数，即 $\boldsymbol{\theta}$ 的分量个数 \\
$\mathbf{w},\mathbf{w}_t$ & 近似价值函数所基于的 $d$ 维权重向量 \\
$w_i,w_{t,i}$ & 可学习权重向量的第 $i$ 个分量 \\
$\hat v(s,\mathbf{w})$ & 给定权重向量 $\mathbf{w}$ 时，状态 $s$ 的近似价值 \\
$v_{\mathbf{w}}(s)$ & $\hat v(s,\mathbf{w})$ 的另一种记法 \\
$\hat q(s,a,\mathbf{w})$ & 给定权重向量 $\mathbf{w}$ 时，状态--动作对 $s,a$ 的近似价值 \\
$\nabla\hat v(s,\mathbf{w})$ & $\hat v(s,\mathbf{w})$ 关于 $\mathbf{w}$ 的偏导数组成的列向量 \\
$\nabla\hat q(s,a,\mathbf{w})$ & $\hat q(s,a,\mathbf{w})$ 关于 $\mathbf{w}$ 的偏导数组成的列向量 \\
$\mathbf{x}(s)$ & 处于状态 $s$ 时可见的特征向量 \\
$\mathbf{x}(s,a)$ & 在状态 $s$ 采取动作 $a$ 时可见的特征向量 \\
$x_i(s),x_i(s,a)$ & 向量 $\mathbf{x}(s)$ 或 $\mathbf{x}(s,a)$ 的第 $i$ 个分量 \\
$\mathbf{x}_t$ & $\mathbf{x}(S_t)$ 或 $\mathbf{x}(S_t,A_t)$ 的简写 \\
$\mathbf{w}^\top\mathbf{x}$ & 向量内积，$\mathbf{w}^\top\mathbf{x}\doteq\sum_i w_i x_i$；例如 $\hat v(s,\mathbf{w})=\mathbf{w}^\top\mathbf{x}(s)$ \\
$\mathbf{v},\mathbf{v}_t$ & 第二个 $d$ 维权重向量，用来学习 $\mathbf{w}$（第 11 章） \\
$\mathbf{z}_t$ & 时刻 $t$ 的 $d$ 维资格迹向量（第 12 章） \\
$\boldsymbol{\theta},\boldsymbol{\theta}_t$ & 目标策略的参数向量（第 13 章） \\
$\pi(a|s,\boldsymbol{\theta})$ & 给定参数向量 $\boldsymbol{\theta}$ 时，在状态 $s$ 采取动作 $a$ 的概率 \\
$\pi_{\boldsymbol{\theta}}$ & 对应于参数 $\boldsymbol{\theta}$ 的策略 \\
$\nabla\pi(a|s,\boldsymbol{\theta})$ & $\pi(a|s,\boldsymbol{\theta})$ 关于 $\boldsymbol{\theta}$ 的偏导数组成的列向量 \\
$J(\boldsymbol{\theta})$ & 策略 $\pi_{\boldsymbol{\theta}}$ 的性能度量 \\
$\nabla J(\boldsymbol{\theta})$ & $J(\boldsymbol{\theta})$ 关于 $\boldsymbol{\theta}$ 的偏导数组成的列向量 \\
$h(s,a,\boldsymbol{\theta})$ & 基于 $\boldsymbol{\theta}$，在状态 $s$ 中选择动作 $a$ 的偏好 \\
$b(a|s)$ & 在学习目标策略 $\pi$ 的同时用于选择动作的行为策略 \\
$b(s)$ & 策略梯度方法中的基线函数 $b:\mathcal{S}\mapsto\mathbb{R}$ \\
$b$ & MDP 或搜索树的分支因子 \\
$\rho_{t:h}$ & 从时刻 $t$ 到 $h$ 的重要性采样比（第 5.5 节） \\
$\rho_t$ & 仅时刻 $t$ 的重要性采样比，$\rho_t\doteq\rho_{t:t}$ \\
$r(\pi)$ & 策略 $\pi$ 的平均奖励（奖励率）（第 10.3 节） \\
$\bar R_t$ & 时刻 $t$ 对 $r(\pi)$ 的估计 \\
$\mu(s)$ & 状态上的同策略分布（第 9.2 节） \\
$\boldsymbol{\mu}$ & 由所有 $s\in\mathcal{S}$ 的 $\mu(s)$ 组成的 $|\mathcal{S}|$ 维向量 \\
$\|v\|_\mu^2$ & 价值函数 $v$ 的 $\mu$ 加权平方范数，即 $\|v\|_\mu^2\doteq\sum_{s\in\mathcal{S}}\mu(s)v(s)^2$ \\
$\eta(s)$ & 每幕中访问状态 $s$ 的期望次数（第 199 页） \\
$\Pi$ & 价值函数的投影算子（第 268 页） \\
$B_\pi$ & 价值函数的贝尔曼算子（第 11.4 节） \\
$\mathbf{A}$ & $d\times d$ 矩阵 $\mathbf{A}\doteq\mathbb{E}\!\left[\mathbf{x}_t(\mathbf{x}_t-\gamma\mathbf{x}_{t+1})^\top\right]$ \\
$\mathbf{b}$ & $d$ 维向量 $\mathbf{b}\doteq\mathbb{E}[R_{t+1}\mathbf{x}_t]$ \\
$\mathbf{w}_{\mathrm{TD}}$ & TD 固定点 $\mathbf{w}_{\mathrm{TD}}\doteq\mathbf{A}^{-1}\mathbf{b}$（一个 $d$ 维向量，第 9.4 节） \\
$\mathbf{I}$ & 单位矩阵 \\
$\mathbf{P}$ & 策略 $\pi$ 下状态转移概率构成的 $|\mathcal{S}|\times|\mathcal{S}|$ 矩阵 \\
$\mathbf{D}$ & 对角线上为 $\boldsymbol{\mu}$ 的 $|\mathcal{S}|\times|\mathcal{S}|$ 对角矩阵 \\
$\mathbf{X}$ & 以 $\mathbf{x}(s)$ 为行的 $|\mathcal{S}|\times d$ 矩阵 \\
$\bar\delta_{\mathbf{w}}(s)$ & 状态 $s$ 上 $v_{\mathbf{w}}$ 的贝尔曼误差（期望 TD 误差）（第 11.4 节） \\
$\bar{\boldsymbol{\delta}}_{\mathbf{w}},\mathrm{BE}$ & 贝尔曼误差向量，其分量为 $\bar\delta_{\mathbf{w}}(s)$ \\
$\overline{\mathrm{VE}}(\mathbf{w})$ & 均方价值误差，$\overline{\mathrm{VE}}(\mathbf{w})\doteq\|v_{\mathbf{w}}-v_\pi\|_\mu^2$（第 9.2 节） \\
$\overline{\mathrm{BE}}(\mathbf{w})$ & 均方贝尔曼误差，$\overline{\mathrm{BE}}(\mathbf{w})\doteq\|\bar{\boldsymbol{\delta}}_{\mathbf{w}}\|_\mu^2$ \\
$\overline{\mathrm{PBE}}(\mathbf{w})$ & 均方投影贝尔曼误差，$\overline{\mathrm{PBE}}(\mathbf{w})\doteq\|\Pi\bar{\boldsymbol{\delta}}_{\mathbf{w}}\|_\mu^2$ \\
$\overline{\mathrm{TDE}}(\mathbf{w})$ & 均方时序差分误差，$\overline{\mathrm{TDE}}(\mathbf{w})\doteq\mathbb{E}_b[\rho_t\delta_t^2]$（第 11.5 节） \\
$\overline{\mathrm{RE}}(\mathbf{w})$ & 均方回报误差（第 11.6 节） \\
\end{longtable}

\section*{第 1 章\quad 引言}

当我们思考学习的本质时，通过与环境相互作用而学习这个想法，很可能是最先浮现在我们脑海中的想法。婴儿玩耍、挥动手臂或四处张望时，并没有明确的教师，但它确实与环境有直接的感觉运动连接。运用这种连接会产生大量关于因果、关于动作后果以及关于为了达到目标应该做什么的信息。在我们的一生中，这类相互作用无疑是我们关于环境和自身知识的重要来源。无论我们是在学习开车，还是在进行一次谈话，我们都会敏锐地意识到环境如何回应我们的所作所为，并试图通过自己的行为影响接下来发生的事情。从相互作用中学习，是几乎所有学习与智能理论背后的基础性思想。

在本书中，我们探讨一种从相互作用中学习的计算方法。我们并不直接从理论上说明人或动物如何学习，而是主要考察理想化的学习情境，并评价各种学习方法的有效性。\footnote{与心理学和神经科学的关系将在第 14 章和第 15 章中概述。}也就是说，我们采取人工智能研究者或工程师的视角。我们探索能够有效解决具有科学或经济意义的学习问题的机器设计，并通过数学分析或计算实验来评价这些设计。我们所探讨的方法称为强化学习；与其他机器学习方法相比，它更专注于从相互作用中进行目标导向的学习。

\subsection*{1.1\quad 强化学习}

强化学习所学习的是做什么，也就是如何把情境映射为动作，以便最大化一个数值奖励信号。学习者不会被告知要采取哪些动作，而必须通过尝试来发现哪些动作会带来最多奖励。在最有趣也最具挑战性的情形中，动作不仅可能影响即时奖励，还可能影响下一个情境，并通过它影响所有后续奖励。试错搜索和延迟奖励这两个特征，是强化学习最重要的两个区别性特征。

强化学习像许多名称以“ing”结尾的主题一样，例如机器学习和登山运动，同时是一类问题、一类在这类问题上效果良好的求解方法，以及研究这一问题及其求解方法的领域。用一个名称称呼这三件事很方便，但同时也必须在概念上把三者区分开。特别地，在强化学习中，问题和求解方法之间的区分非常重要；不能做出这种区分，是许多混淆的来源。

我们用动力系统理论中的思想形式化强化学习问题，具体地说，把它看作不完全已知的马尔可夫决策过程的最优控制。这种形式化的细节必须等到第 3 章再说，但基本思想只是要捕捉这样一个真实问题中最重要的方面：一个学习智能体随着时间与其环境相互作用，以实现某个目标。学习智能体必须能够在某种程度上感知其环境的状态，并且必须能够采取影响该状态的动作。智能体还必须有一个或多个与环境状态有关的目标。马尔可夫决策过程旨在以尽可能简单而又不平凡化任何一方面的形式，把这三个方面，即感觉、动作和目标，都包括进来。任何适合求解这类问题的方法，我们都认为是强化学习方法。

强化学习不同于监督学习，后者是当前机器学习领域大多数研究所考察的学习类型。监督学习是从一个由有知识的外部监督者提供的带标签样例训练集中学习。每个样例都是某个情境的描述，并伴随一个说明，也就是系统在该情境下应采取的正确动作的标签；这个动作通常是识别该情境所属的类别。这类学习的目标，是让系统能够外推或泛化它的反应，从而在训练集中没有出现过的情境下也能正确行动。这是一类重要的学习，但仅靠它并不足以从相互作用中学习。在交互式问题中，通常不现实的是，获得既正确又能代表智能体必须行动的所有情境的期望行为样例。在未被探明的领域中，也就是人们最期待学习发挥作用的地方，智能体必须能够从自身经验中学习。

强化学习也不同于机器学习研究者所称的无监督学习；无监督学习通常是要在无标签数据集合中发现隐藏结构。监督学习和无监督学习这两个术语似乎可以穷尽机器学习范式，但事实并非如此。由于强化学习不依赖正确行为的样例，人们也许会倾向于把强化学习看作一种无监督学习；然而强化学习试图最大化奖励信号，而不是试图发现隐藏结构。在智能体经验中揭示结构当然可能对强化学习有用，但它本身并没有处理强化学习中最大化奖励信号的问题。因此，我们把强化学习看作第三种机器学习范式，与监督学习、无监督学习以及也许还有其他范式并列。

强化学习中会出现而其他学习类型中不会出现的一个挑战，是探索与利用之间的权衡。为了获得大量奖励，强化学习智能体必须偏好那些它过去尝试过并发现能有效产生奖励的动作。但是，为了发现这样的动作，它又必须尝试以前没有选择过的动作。智能体必须利用已经经历过的内容以获得奖励，但也必须探索以便将来做出更好的动作选择。困难在于，若只追求探索或只追求利用，都会导致任务失败。智能体必须尝试多种动作，并逐渐偏向那些看起来最好的动作。在随机性任务中，每个动作都必须被尝试多次，才能得到其期望奖励的可靠估计。探索--利用困境已经被数学家密集研究了几十年，却仍未被解决。眼下我们只需注意，在监督学习和无监督学习中，至少在这些范式最纯粹的形式下，平衡探索与利用的整个问题甚至不会出现。

强化学习的另一个关键特征是，它明确考虑一个目标导向智能体与不确定环境相互作用的整体问题。这与许多只考虑子问题、却不说明这些子问题如何放入更大图景中的方法形成对比。例如，我们已经提到，机器学习研究很大一部分关注监督学习，却没有明确说明这种能力最终将如何发挥作用。另一些研究者发展了带有一般目标的规划理论，却没有考虑规划在实时决策中的作用，也没有考虑规划所需预测模型来自何处。尽管这些方法产生了许多有用结果，但它们集中于孤立子问题这一点是一个显著限制。

强化学习采取相反的路线，从一个完整的、交互的、追求目标的智能体出发。所有强化学习智能体都有明确目标，能够感知其环境的若干方面，并能够选择动作来影响其环境。此外，通常从一开始就假定，智能体必须在面对环境的显著不确定性时仍然运作。当强化学习涉及规划时，它必须处理规划与实时动作选择之间的相互作用，也必须处理环境模型如何获得和改进的问题。当强化学习涉及监督学习时，它这样做是出于特定原因，这些原因决定了哪些能力是关键的，哪些不是。为了推动学习研究取得进展，重要的子问题必须被分离出来加以研究，但它们应当是在完整的、交互的、追求目标的智能体中起明确作用的子问题，即使完整智能体的全部细节尚不能被填补。

所谓完整的、交互的、追求目标的智能体，并不总是指完整的生物体或机器人。这些当然是例子，但一个完整的、交互的、追求目标的智能体也可以是更大行为系统的一个组成部分。在这种情况下，智能体直接与更大系统的其余部分相互作用，并间接与更大系统的环境相互作用。一个简单例子是，某个智能体监测机器人电池的电量，并向机器人的控制架构发送命令。这个智能体的环境是机器人其余部分以及机器人的环境。人们必须超越最明显的智能体及其环境的例子，才能体会强化学习框架的一般性。

现代强化学习最令人兴奋的方面之一，是它与其他工程和科学学科之间实质而富有成果的相互作用。强化学习是人工智能和机器学习中一项持续数十年的趋势的一部分：它们正走向与统计、优化以及其他数学主题的更大整合。例如，某些强化学习方法使用参数化近似器学习的能力，处理了运筹学和控制理论中经典的“维数灾难”。更有特色的是，强化学习也与心理学和神经科学发生了强烈互动，并且双方都获得了实质收益。在所有机器学习形式中，强化学习最接近人类和其他动物所做的那种学习，而且强化学习的许多核心算法最初受到生物学习系统的启发。强化学习也有所回馈：一方面通过一个更好匹配某些经验数据的动物学习心理学模型，另一方面通过一个关于大脑奖励系统若干部分的有影响力模型。本书主体发展的是与工程和人工智能有关的强化学习思想，而与心理学和神经科学的联系将在第 14 章和第 15 章中概述。

最后，强化学习也是人工智能中一个更大趋势的一部分，即重新回到简单的一般原则。自 20 世纪 60 年代后期以来，许多人工智能研究者推定不存在有待发现的一般原则，智能反而源自拥有大量特殊用途的技巧、过程和启发式规则。有时人们会说，如果我们能把足够多的相关事实输入机器，比如一百万、一千万或十亿个事实，那么机器就会变得有智能。基于一般原则的方法，例如搜索或学习，被称为“弱方法”，而基于特定知识的方法被称为“强方法”。这种观点今天仍然常见，但已不占主导地位。从我们的角度看，它只是为时过早：人们在寻找一般原则上投入的努力太少，不足以得出没有一般原则的结论。现代人工智能现在包含大量寻找学习、搜索和决策制定的一般原则的研究。钟摆会摆回多远尚不清楚，但强化学习研究无疑是人工智能重新转向更简单、更少的一般原则这一摆动的一部分。

\subsection*{1.2\quad 例子}

理解强化学习的一种好方法，是考察一些指导其发展的例子和可能应用。

\begin{itemize}
\item 一位国际象棋大师走了一步棋。这个选择既受到规划的影响，即预期可能的应招和反应招，也受到对特定局面和走法可取性的即时、直觉判断的影响。

\item 一个自适应控制器实时调整炼油厂运行的参数。这个控制器根据指定的边际成本优化产量、成本和质量之间的权衡，而不是严格固守工程师最初建议的设定点。

\item 一只瞪羚幼崽出生后几分钟就挣扎着站起来。半小时后，它已经能以每小时 20 英里的速度奔跑。

\item 一个移动机器人决定它是否应当进入一个新房间寻找更多垃圾来收集，还是开始设法返回电池充电站。它根据电池当前电量，以及过去能以多快、多容易找到充电器，来做出这个决定。

\item Phil 准备他的早餐。仔细考察后，即使这种看似平凡的活动，也揭示出一张复杂的条件行为网络以及相互嵌套的目标--子目标关系：走到橱柜前，打开橱柜，选择一盒麦片，然后伸手、抓住并取回盒子。为了取得碗、勺子和牛奶盒，还需要其他复杂的、调谐过的交互式行为序列。每一步都涉及一系列眼动，以获取信息并引导伸手和移动。关于如何搬运这些物体，或者在取得其他物体前是否最好先把其中一些送到餐桌上，人会不断做出快速判断。每一步都受目标引导，例如抓住勺子或到冰箱前；同时它也服务于其他目标，例如拿到勺子，以便在麦片准备好之后用它吃饭，并最终获得营养。无论 Phil 是否意识到这一点，他都在获取关于自己身体状态的信息，这些信息决定了他的营养需求、饥饿程度和食物偏好。
\end{itemize}

这些例子共享一些非常基本、以至于容易被忽视的特征。所有例子都涉及一个主动决策智能体与其环境之间的相互作用；在这种相互作用中，智能体尽管面对环境的不确定性，仍试图实现一个目标。智能体的动作被允许影响环境的未来状态，例如下一个棋局、炼油厂储罐的液位、机器人的下一个位置以及其电池未来电量，从而影响以后可供智能体使用的动作和机会。正确选择需要考虑动作间接的、延迟的后果，因此可能需要远见或规划。

与此同时，在所有这些例子中，动作的效果都不能被完全预测；因此，智能体必须频繁监测其环境并做出适当反应。例如，Phil 必须看着他倒进麦片碗里的牛奶，以免牛奶溢出。所有这些例子都涉及明确的目标，所谓明确，是指智能体能够根据它能直接感知的内容，判断朝向目标的进展。棋手知道自己是否获胜，炼油厂控制器知道正在生产多少石油，瞪羚幼崽知道自己何时摔倒，移动机器人知道自己的电池何时电量耗尽，而 Phil 知道自己是否正在享用早餐。

在所有这些例子中，智能体都能用它的经验随时间改善表现。棋手改进他用来评价局面的直觉，从而改善棋艺；瞪羚幼崽提高它奔跑的效率；Phil 学会简化制作早餐的过程。智能体在任务开始时带来的知识，或来自相关任务的先前经验，或由设计、演化内置其中，会影响什么是有用或容易学习的；但与环境的相互作用对于调整行为、以利用任务的特定特征来说是必不可少的。

\subsection*{1.3\quad 强化学习的要素}

除了智能体和环境之外，还可以识别出强化学习系统的四个主要子要素：策略、奖励信号、价值函数，以及可选的环境模型。

策略定义了学习智能体在给定时刻的行为方式。粗略地说，策略是从被感知到的环境状态到在那些状态下应采取动作的映射。它对应于心理学中所谓的一组刺激--反应规则或联结。在某些情况下，策略可以是一个简单函数或查找表；而在另一些情况下，它可能涉及大量计算，例如搜索过程。策略是强化学习智能体的核心，因为仅凭策略本身就足以决定行为。一般而言，策略可以是随机性的，为每个动作指定概率。

奖励信号定义了强化学习问题的目标。在每个时间步，环境向强化学习智能体发送一个称为奖励的单个数字。智能体唯一的目标，是最大化它长期接收的总奖励。因此，奖励信号为智能体定义了什么是好事件、什么是坏事件。在生物系统中，我们可以把奖励看作类似于愉悦或痛苦的经验。它们是智能体所面对问题的即时而定义性的特征。奖励信号是改变策略的主要依据；如果策略选择的某个动作随后得到低奖励，那么将来在该情境下，策略可能会被改变为选择其他动作。一般而言，奖励信号可以是环境状态和所采取动作的随机函数。

奖励信号指出在即时意义上什么是好的，而价值函数则规定从长远看什么是好的。粗略地说，一个状态的价值，是智能体从该状态开始，在未来可以预期累积的总奖励量。奖励决定环境状态即时的、内在的可取性，而价值在考虑可能随之而来的状态以及那些状态中可获得的奖励之后，指出状态的长期可取性。例如，一个状态可能总是产生较低的即时奖励，但仍然具有高价值，因为它经常后接其他会产生高奖励的状态。反过来也可能成立。用人的类比来说，奖励有些像愉悦（若高）和痛苦（若低），而价值对应于一种更精细、更有远见的判断：我们对于环境处在某个特定状态有多满意或不满意。

从某种意义上说，奖励是首要的，而作为奖励预测的价值是次要的。没有奖励就不会有价值，估计价值的唯一目的也是获得更多奖励。尽管如此，在做出和评价决策时，我们最关心的却是价值。动作选择是根据价值判断做出的。我们寻找能产生最高价值状态的动作，而不是产生最高奖励的动作，因为这些动作从长期看为我们获得最多奖励。不幸的是，确定价值比确定奖励困难得多。奖励基本上由环境直接给出，但价值必须根据智能体一生中产生的观察序列进行估计并反复重新估计。事实上，我们所考虑的几乎所有强化学习算法中，最重要的组成部分都是一个有效估计价值的方法。价值估计的核心作用，可以说是过去六十年中关于强化学习所学到的最重要事情。

某些强化学习系统的第四个、也是最后一个要素，是环境模型。这是一种模仿环境行为的东西，或者更一般地说，它允许我们对环境将如何行为做出推断。例如，给定一个状态和动作，模型可能预测随后的下一个状态和下一个奖励。模型用于规划；所谓规划，是指在实际经历可能的未来情境之前，先通过考虑它们来决定行动路线的任何方式。使用模型和规划来求解强化学习问题的方法称为基于模型的方法，与之相对的是更简单的无模型方法；后者是明确的试错学习者，几乎可以看作规划的反面。在第 8 章中，我们将考察同时通过试错学习、学习环境模型，并使用该模型进行规划的强化学习系统。现代强化学习覆盖了从低层次试错学习到高层次深思熟虑式规划的整个谱系。

\subsection*{1.4\quad 局限与范围}

强化学习高度依赖状态这一概念；状态既作为策略和价值函数的输入，也作为模型的输入和输出。非正式地说，我们可以把状态看作传递给智能体的一种信号，它在某个特定时刻表达“环境如何”的某种感觉。我们在这里使用的状态的形式化定义，由第 3 章介绍的马尔可夫决策过程框架给出。不过，更一般地说，我们鼓励读者遵循这种非正式含义，把状态看作智能体关于其环境可获得的任何信息。实际上，我们假定状态信号由某个预处理系统产生，而这个系统名义上是智能体环境的一部分。本书不处理构造、改变或学习状态信号的问题（除第 17.3 节简要涉及外）。我们采取这种做法，并不是因为认为状态表示不重要，而是为了完全集中于决策问题。换言之，本书关心的不是设计状态信号，而是作为可用状态信号的函数来决定采取什么动作。

本书中考虑的大多数强化学习方法，都围绕估计价值函数来组织，但求解强化学习问题并不严格要求这样做。例如，遗传算法、遗传规划、模拟退火以及其他优化方法等求解方法，从不估计价值函数。这些方法会在较长时间内与环境的一个单独实例交互，应用多种静态策略。获得最多奖励的策略以及它们的随机变体，会被带入下一代策略，这一过程不断重复。我们称这些为演化方法，因为它们的运作方式类似于生物进化产生具有熟练行为的有机体，即使这些有机体在各自的一生中并不学习。如果策略空间足够小，或者能够被结构化，使得好策略常见或容易找到，或者有大量时间可用于搜索，那么演化方法可能是有效的。此外，在学习智能体不能感知其环境完整状态的问题上，演化方法具有优势。

我们的关注点是那些在与环境相互作用时学习的强化学习方法，而演化方法并不这样做。能够利用个体行为交互细节的方法，在许多情况下可能比演化方法高效得多。演化方法忽略了强化学习问题中大量有用结构：它们没有利用正在搜索的策略是从状态到动作的函数这一事实；它们不注意个体一生中经过了哪些状态，或选择了哪些动作。在某些情况下，这些信息可能具有误导性（例如状态被错误感知时），但更多时候它应当能够使搜索更高效。尽管进化和学习共享许多特征，并自然地协同工作，我们并不认为单独的演化方法特别适合强化学习问题；因此，本书不讨论它们。

\subsection*{1.5\quad 一个扩展示例：井字棋}

为了说明强化学习的一般思想并把它与其他方法对比，我们接下来更详细地考察一个单一例子。

考虑熟悉的儿童游戏井字棋。两名玩家轮流在一个 $3\times3$ 棋盘上下棋。一名玩家下 X，另一名玩家下 O，直到某名玩家通过横向、纵向或对角线放置三个标记连成一线而获胜；右侧所示游戏中，X 玩家就是这样获胜的。如果棋盘填满而双方都没有三子连线，则游戏为平局。由于有技巧的玩家可以做到永不失败，我们假设自己面对的是一个不完美的玩家，也就是它有时会下错，从而允许我们获胜的玩家。实际上，眼下我们把平局和失败都看作对我们同样糟糕。我们如何构造一个玩家，使它能发现对手下法中的不完善之处，并学会最大化自己的获胜机会呢？

\begin{figure}[htbp]
  \centering
  \includegraphics[width=0.26\textwidth]{tex/强化学习/figures/tic-tac-toe-board.png}
\end{figure}

尽管这是一个简单问题，它却不容易用经典技术令人满意地解决。例如，博弈论中的经典“极小化极大”解法在这里并不正确，因为它假定对手采用某种特定下法。例如，极小化极大玩家永远不会到达一个它可能输掉的棋局状态，即使事实上由于对手下错，它总是能从那个状态获胜。用于序贯决策问题的经典优化方法，例如动态规划，可以为任意对手计算一个最优解，但要求把该对手的完整规格作为输入，其中包括对手在每个棋盘状态下做出每一步的概率。我们假设对于这个问题，这些信息并不是先验可得的；对于绝大多数具有实际意义的问题也是如此。另一方面，这类信息可以从经验中估计；在本例中，就是通过与对手下许多盘棋来估计。针对这个问题，人们差不多能做的最好事情，是首先学习一个关于对手行为的模型，达到某种置信水平，然后应用动态规划，根据近似的对手模型计算最优解。归根到底，这与本书后面将考察的某些强化学习方法并没有那么不同。

应用于这个问题的演化方法，会直接搜索可能策略的空间，以找到一个有很高概率击败对手的策略。在这里，策略是一条规则，它告诉玩家在游戏的每个状态下该走哪一步，也就是在三乘三棋盘上 X 和 O 的每一种可能配置下该走哪一步。对于每个被考虑的策略，其获胜概率的估计会通过与对手下若干盘棋而获得。这个评价随后会指导接下来考虑哪个或哪些策略。典型的演化方法会在策略空间中爬山，不断生成并评价策略，试图获得增量式改进。或者，也许可以使用一种遗传风格的算法，维护并评价一个策略种群。实际上，可以应用数以百计的不同优化方法。

下面说明如何用一种使用价值函数的方法来处理井字棋问题。首先，我们会建立一张数字表，游戏的每个可能状态对应一个数字。每个数字都是从该状态开始我方获胜概率的最新估计。我们把这个估计看作该状态的价值，整张表就是学到的价值函数。如果从状态 A 获胜的当前概率估计高于从状态 B 获胜的估计，那么状态 A 的价值就高于状态 B，或者说状态 A 被认为比状态 B “更好”。假设我们总是下 X，那么所有已有三个 X 连成一线的状态，其获胜概率为 $1$，因为我们已经获胜。同样，所有已有三个 O 连成一线的状态，或者棋盘已满的状态，其正确概率为 $0$，因为我们无法从这些状态获胜。我们把所有其他状态的初始价值设为 $0.5$，表示我们猜测有 $50\%$ 的获胜机会。

然后，我们与对手下许多盘棋。为了选择自己的走法，我们考察每个可能走法会导致的状态（棋盘上每个空格对应一个），并在表中查找它们的当前价值。大多数时候我们贪心地走棋，选择导致最大价值状态的走法，也就是导致最高估计获胜概率的走法。不过，我们偶尔会从其他走法中随机选择。这些称为探索性走法，因为它们使我们经历一些否则可能永远看不到的状态。一盘棋中实际走过和考虑过的走法序列，可以像图 1.1 那样画出来。

在下棋过程中，我们改变自己在棋局中所处状态的价值。我们试图使它们成为获胜概率的更准确估计。为此，正如图 1.1 中的箭头所暗示的，我们把每次贪心走法之后状态的价值“回传”到走法之前的状态。更准确地说，较早状态的当前价值被更新为更接近较晚状态的价值。可以通过把较早状态的价值沿着通向较晚状态价值的方向移动一小部分来完成这一点。如果用 $S_t$ 表示贪心走法之前的状态，用 $S_{t+1}$ 表示走法之后的状态，那么对 $S_t$ 的估计价值（记为 $V(S_t)$）的更新可以写成
\[
  V(S_t)\leftarrow V(S_t)+\alpha\bigl[V(S_{t+1})-V(S_t)\bigr],
\]
其中 $\alpha$ 是一个称为步长参数的小正分数，它影响学习速率。这个更新规则是时序差分学习方法的一个例子；之所以这样称呼，是因为它的改变基于相邻两个时刻估计值之间的差 $V(S_{t+1})-V(S_t)$。

\begin{figure}[htbp]
  \centering
  \includegraphics[width=0.62\textwidth]{tex/强化学习/figures/figure-1-1.png}
  \caption{井字棋走法序列。黑色实线表示一盘棋中实际采取的走法；虚线表示我们（强化学习玩家）考虑过但没有采取的走法。我们的第二步是一次探索性走法，这意味着尽管另一个兄弟走法，即通向 $e^*$ 的走法，排名更高，我们仍采取了这一步。探索性走法不会导致任何学习，但我们的其他每一步都会导致学习；正如红色箭头所示，估计价值会按照正文详述，从树中较晚的节点向较早的节点上移。}
\end{figure}

上面描述的方法在这个任务上表现相当好。例如，如果步长参数随时间适当降低，那么对于任意固定对手，该方法都会收敛到在我方最优下法下，从每个状态获胜的真实概率。此外，随后采取的走法（探索性走法除外）事实上就是对抗这个不完美对手的最优走法。换言之，该方法会收敛到一个与这个对手下棋的最优策略。如果步长参数随时间并不一直降低到零，那么这个玩家在面对下法缓慢变化的对手时也能表现良好。

这个例子说明了演化方法与学习价值函数的方法之间的差异。为了评价一个策略，演化方法保持该策略固定，并与对手下许多盘棋，或使用对手模型模拟许多盘棋。获胜频率给出了采用该策略时获胜概率的无偏估计，并可用来指导下一次策略选择。但是，每一次策略改变都只在许多盘棋之后做出，并且每盘棋只使用最终结果：棋局进行过程中发生了什么被忽略。例如，如果玩家获胜，那么它在这盘棋中的所有行为都会获得功劳，而不管哪些具体走法可能对获胜至关重要。甚至从未发生过的走法也会获得功劳！相比之下，价值函数方法允许评价个体状态。最后，演化方法和价值函数方法都搜索策略空间，但学习价值函数利用了下棋过程中可获得的信息。

这个简单例子说明了强化学习方法的一些关键特征。第一，它强调在与环境相互作用时学习；本例中，环境是对手玩家。第二，它有明确目标，而正确行为需要考虑选择的延迟影响，因此需要规划或远见。例如，这个简单的强化学习玩家会学会为短视的对手设置多步陷阱。强化学习解法的一个显著特征是，它可以在不使用对手模型、也不显式搜索未来状态和动作可能序列的情况下，实现规划和前瞻的效果。

虽然这个例子说明了强化学习的一些关键特征，但它太简单了，可能会让人误以为强化学习比实际更受限制。尽管井字棋是一个双人游戏，强化学习也适用于不存在外部对手的情形，也就是“与自然博弈”的情形。强化学习也不限于那些行为分解为独立幕的问题，例如每盘井字棋只在每幕结束时给出奖励的情况。当行为无限持续、且可以在任意时刻接收不同大小的奖励时，它同样适用。强化学习也适用于那些甚至不能像井字棋走法那样分解为离散时间步的问题。一般原则同样适用于连续时间问题，不过理论会更加复杂，我们在这份导论性处理里略去它。

井字棋有一个相对较小的有限状态集，而强化学习可以在状态集非常大、甚至无限的情况下使用。例如，Gerry Tesauro（1992, 1995）把上面描述的算法与人工神经网络结合起来，学习下双陆棋；双陆棋大约有 $10^{20}$ 个状态。状态如此之多，以至于不可能经历超过其中很小一部分的状态。Tesauro 的程序学会了下得远好于此前任何程序，并最终好于世界上最强的人类玩家（第 16.1 节）。人工神经网络为程序提供了从经验中泛化的能力，使它能在新状态中，根据其网络从过去遇到的相似状态中保存的信息来选择走法。强化学习系统在具有如此大状态集的问题上能工作得多好，与它从过去经验中泛化得有多恰当密切相关。正是在这个角色上，我们最需要把监督学习方法与强化学习结合起来。人工神经网络和深度学习（第 9.6 节）并不是做到这一点的唯一方法，也未必是最好的方法。

在这个井字棋例子中，学习开始时除了游戏规则之外没有先验知识，但强化学习绝不意味着对学习和智能采取一块白板式的看法。相反，先验信息可以通过多种方式并入强化学习，而这对于高效学习可能至关重要（例如见第 9.5、17.4 和 13.1 节）。在井字棋例子中，我们还能够访问真实状态；而强化学习也可以在部分状态隐藏，或不同状态在学习者看来相同的情况下应用。

最后，井字棋玩家能够前瞻，并知道它每个可能走法会导致哪些状态。为此，它必须拥有一个游戏模型，使它能够预见环境会如何响应它也许从未采取过的走法而改变。许多问题都像这样，但在另一些问题中，甚至缺少动作效果的短期模型。无论哪种情况，强化学习都可以应用。模型不是必需的，但如果模型可用或可学习，就可以很容易地加以利用（第 8 章）。

另一方面，也存在完全不需要任何环境模型的强化学习方法。无模型系统甚至不能思考环境会如何响应单个动作而改变。从这个意义上说，井字棋玩家相对于它的对手是无模型的：它没有任何形式的对手模型。由于模型必须相当准确才有用，当求解问题的真正瓶颈在于构造足够准确的环境模型时，无模型方法可能比更复杂的方法更有优势。无模型方法也是基于模型方法的重要构件。本书会先用几章讨论无模型方法，然后再讨论它们如何作为更复杂的基于模型方法的组成部分使用。

强化学习可以在一个系统的高层次和低层次上使用。虽然井字棋玩家只学习游戏的基本走法，但没有什么阻止强化学习在更高层次上工作；在那里，每个“动作”本身都可能是应用某种相当复杂的问题求解方法。在分层学习系统中，强化学习可以同时在多个层次上工作。

\noindent\textbf{练习 1.1：自我博弈}\quad 假设上面描述的强化学习算法不是与随机对手对弈，而是与自身对弈，双方都在学习。你认为这种情况下会发生什么？它会学到不同的走法选择策略吗？

\noindent\textbf{练习 1.2：对称性}\quad 许多井字棋局面看起来不同，但由于对称性实际上相同。我们该如何修改上面描述的学习过程，以利用这一点？这会以哪些方式改善学习过程？再想一想。假设对手没有利用对称性。在这种情况下，我们应当利用吗？那么，对称等价的位置是否必然应具有相同价值？

\noindent\textbf{练习 1.3：贪心下法}\quad 假设强化学习玩家是贪心的，也就是说，它总是下到自己评价最高的位置。它可能会比非贪心玩家学得更好，还是更差？可能会出现什么问题？

\noindent\textbf{练习 1.4：从探索中学习}\quad 假设学习更新发生在所有走法之后，包括探索性走法。如果步长参数随时间适当降低（但探索倾向不降低），那么状态价值会收敛到另一组概率。当我们从探索性走法中学习和不从探索性走法中学习时，从概念上说，分别计算的是哪两组概率？假设我们确实持续进行探索性走法，哪一组概率可能更值得学习？哪一种会带来更多胜利？

\noindent\textbf{练习 1.5：其他改进}\quad 你能想到其他改进强化学习玩家的方法吗？你能想到解决这里提出的井字棋问题的任何更好方法吗？

\subsection*{1.6\quad 小结}

强化学习是一种计算方法，用来理解并自动化目标导向的学习和决策制定。它区别于其他计算方法之处在于，它强调智能体通过与环境直接相互作用来学习，而不要求范例式监督或完整的环境模型。在我们看来，强化学习是第一个认真处理如下计算问题的领域：为了实现长期目标而与环境相互作用时，学习会产生哪些计算问题。

强化学习使用马尔可夫决策过程的形式化框架，以状态、动作和奖励来定义学习智能体与其环境之间的相互作用。这个框架旨在用一种简单方式表示人工智能问题的基本特征。这些特征包括因果感、不确定性和非确定性感，以及显式目标的存在。

价值和价值函数的概念，是本书所考虑的大多数强化学习方法的关键。我们的立场是，价值函数对于在策略空间中高效搜索很重要。价值函数的使用，使强化学习方法区别于演化方法；后者由对整个策略的评价来引导，直接在策略空间中搜索。

\subsection*{1.7\quad 强化学习的早期历史}

强化学习的早期历史有两条主要线索，二者都漫长而丰富，在现代强化学习中交织之前曾被独立追求。一条线索涉及试错学习，起源于动物学习心理学。这条线索贯穿了人工智能中一些最早期的工作，并在 20 世纪 80 年代初导致强化学习的复兴。另一条线索涉及最优控制问题，以及使用价值函数和动态规划求解它的方法。在很大程度上，这条线索并不涉及学习。两条线索大多彼此独立，但在第三条较不清晰的线索周围发生了某种程度的关联；这第三条线索涉及时序差分方法，例如本章井字棋例子中使用的方法。三条线索在 20 世纪 80 年代末汇合，产生了我们在本书中呈现的现代强化学习领域。

我们最熟悉、在这段简史中也最有话可说的是聚焦于试错学习的线索。不过在进入它之前，我们先简要讨论最优控制这条线索。

“最优控制”这一术语在 20 世纪 50 年代后期开始使用，用来描述设计控制器的问题；该控制器要随着时间最小化或最大化对动力系统行为的某个度量。Richard Bellman 等人在 20 世纪 50 年代中期发展出一种处理这一问题的方法，它扩展了 Hamilton 和 Jacobi 在 19 世纪的理论。这种方法使用动力系统状态和价值函数，或“最优回报函数”的概念，来定义一个泛函方程；这个方程现在通常称为贝尔曼方程。通过求解这一方程来解决最优控制问题的一类方法，被称为动态规划（Bellman, 1957a）。Bellman（1957b）还引入了最优控制问题的离散随机版本，即马尔可夫决策过程（MDP）。Ronald Howard（1960）为 MDP 设计了策略迭代方法。所有这些都是现代强化学习理论和算法背后的基本要素。

动态规划被广泛认为是求解一般随机最优控制问题的唯一可行方式。它遭受 Bellman 所称的“维数灾难”，意思是其计算需求随状态变量数量呈指数增长；但它仍然远比任何其他一般方法更高效、更广泛适用。自 20 世纪 50 年代后期以来，动态规划已经得到广泛发展，包括对部分可观测 MDP 的扩展（Lovejoy, 1991 综述）、许多应用（White, 1985, 1988, 1993 综述）、近似方法（Rust, 1996 综述）以及异步方法（Bertsekas, 1982, 1983）。动态规划有许多出色的现代处理方式（例如 Bertsekas, 2005, 2012；Puterman, 1994；Ross, 1983；Whittle, 1982, 1983）。Bryson（1996）提供了一部权威的最优控制史。

一方面是最优控制和动态规划，另一方面是学习，它们之间的联系很慢才被认识到。我们无法确定造成这种分离的原因，但其主要原因很可能是所涉学科之间的分离以及目标不同。另一个促成因素可能是，人们普遍把动态规划看作一种离线计算，它本质上依赖准确的系统模型和贝尔曼方程的解析解。此外，动态规划最简单的形式是一种沿时间向后进行的计算，因此很难看出它如何参与一个必须向前进行的学习过程。一些最早的动态规划工作，例如 Bellman 和 Dreyfus（1959）的工作，现在可能会被归类为遵循一种学习方法。Witten（1977）的工作（下文讨论）当然有资格被看作学习与动态规划思想的结合。Werbos（1987）明确主张动态规划与学习方法之间应有更大关联，并主张动态规划与理解神经和认知机制的相关性。至于动态规划方法与在线学习的完全整合，则直到 Chris Watkins 1989 年的工作才出现；他用 MDP 形式主义处理强化学习的方式已被广泛采用。此后，这些关系已被许多研究者广泛发展，尤其是 Dimitri Bertsekas 和 John Tsitsiklis（1996）；他们创造了“神经动态规划”一词，用来指动态规划和人工神经网络的结合。当前使用的另一个术语是“近似动态规划”。这些不同方法强调主题的不同方面，但它们都与强化学习共享一种兴趣：绕开动态规划的经典缺点。

在某种意义上，我们也认为最优控制中的全部工作都属于强化学习工作。我们把强化学习方法定义为求解强化学习问题的任何有效方式；现在已经清楚，这些问题与最优控制问题密切相关，特别是与那些被表述为 MDP 的随机最优控制问题密切相关。因此，我们也必须把最优控制的求解方法，例如动态规划，视为强化学习方法。由于几乎所有常规方法都要求完全知道要控制的系统，说它们是强化学习的一部分会感觉有点不自然。另一方面，许多动态规划算法是增量式和迭代式的。像学习方法一样，它们通过连续近似逐渐达到正确答案。正如本书其余部分将展示的，这些相似性远不止表面。完整知识和不完整知识情形下的理论和求解方法关系如此密切，以至于我们觉得必须把它们作为同一主题的一部分一起考虑。

现在让我们回到通向现代强化学习领域的另一条主要线索，也就是以试错学习思想为中心的线索。这里我们只触及主要联系点，并将在第 14.3 节更详细地讨论这一主题。按照美国心理学家 R. S. Woodworth（1938）的说法，试错学习这一思想至少可以追溯到 19 世纪 50 年代 Alexander Bain 关于通过“摸索和实验”学习的讨论，更明确地说，也可以追溯到英国动物行为学家兼心理学家 Conway Lloyd Morgan 于 1894 年使用这一术语描述其动物行为观察。也许第一个把试错学习的本质简洁表达为学习原则的人，是 Edward Thorndike：

\begin{quote}
在同一情境下做出的若干反应中，在其他条件相同的情况下，那些伴随或紧接着令动物满意的反应，会与该情境更牢固地联结起来；因此，当该情境再次出现时，它们更可能再次发生。那些伴随或紧接着令动物不适的反应，在其他条件相同的情况下，会削弱它们与该情境的联结；因此，当该情境再次出现时，它们较不可能发生。满意或不适越强，联结的增强或削弱就越强。（Thorndike, 1911, p. 244）
\end{quote}

Thorndike 称之为“效果律”，因为它描述了强化事件对选择动作倾向的影响。Thorndike 后来修改了该定律，以解释后来关于动物学习的数据（例如奖励和惩罚效果之间的差异）；该定律的各种形式在学习理论家之间引发了相当多争议（例如见 Gallistel, 2005；Herrnstein, 1970；Kimble, 1961, 1967；Mazur, 1994）。尽管如此，效果律以某种形式被广泛视为许多行为背后的基本原则（例如 Hilgard and Bower, 1975；Dennett, 1978；Campbell, 1960；Cziko, 1995）。它是 Clark Hull（1943, 1952）有影响力的学习理论，以及 B. F. Skinner（1938）有影响力的实验方法的基础。

在动物学习语境中，“强化”这一术语是在 Thorndike 表述效果律之后很久才开始使用的。据我们所知，它最早出现在这一语境中，是在 1927 年 Pavlov 关于条件反射专著的英文译本中。Pavlov 把强化描述为：由于动物在与另一个刺激或某个反应具有适当时间关系的情况下接收了一个刺激，即强化物，从而使某种行为模式得到加强。一些心理学家把强化的思想扩展为既包括行为的弱化也包括行为的强化，并把强化物的思想扩展为也可能包括刺激的省略或终止。要被视为强化物，强化或弱化必须在强化物撤除后仍持续；一个只是吸引动物注意，或只是激发其行为而不产生持久变化的刺激，不会被视为强化物。

在计算机中实现试错学习的想法，出现在关于人工智能可能性的最早思考之中。在 1948 年的一份报告中，Alan Turing 描述了一种“愉悦--痛苦系统”的设计，其运作方式与效果律一致：

\begin{quote}
当达到一个动作尚未确定的配置时，会对缺失数据做出随机选择，并把相应条目暂时写入描述中，然后加以应用。当痛苦刺激出现时，所有暂时条目都被取消；当愉悦刺激出现时，它们全部变为永久条目。（Turing, 1948）
\end{quote}

许多巧妙的机电机器被构造出来，用以展示试错学习。最早的可能是 Thomas Ross（1933）建造的一台机器，它能够通过一个简单迷宫找到路径，并通过开关设置记住这条路径。1951 年，W. Grey Walter 建造了他的“机械龟”（Walter, 1950）的一个版本，能够进行一种简单形式的学习。1952 年，Claude Shannon 展示了一只名为 Theseus 的走迷宫老鼠，它使用试错法在迷宫中找到路径，而迷宫本身通过地板下的磁铁和继电器记住成功方向（另见 Shannon, 1951）。J. A. Deutsch（1954）描述了一台基于其行为理论（Deutsch, 1953）的迷宫求解机器，它具有某些与基于模型的强化学习（第 8 章）相同的性质。Marvin Minsky（1954）在其博士论文中讨论了强化学习的计算模型，并描述了他建造的一台模拟机器；该机器由他称为 SNARCs（Stochastic Neural-Analog Reinforcement Calculators，随机神经模拟强化计算器）的部件组成，意在类似于大脑中可修改的突触连接（第 15 章）。网站 cyberneticzoo.com 包含关于这些以及许多其他机电学习机器的大量信息。

构造机电学习机器让位于为数字计算机编程，使其执行各种类型的学习，其中一些实现了试错学习。Farley 和 Clark（1954）描述了一种神经网络学习机器的数字模拟，它通过试错学习。但他们的兴趣很快从试错学习转向泛化和模式识别，也就是从强化学习转向监督学习（Clark and Farley, 1955）。这开启了一种关于这些学习类型之间关系的混淆模式。许多研究者似乎认为自己在研究强化学习，而实际上他们研究的是监督学习。例如，Rosenblatt（1962）以及 Widrow 和 Hoff（1960）等人工神经网络先驱显然受到强化学习的激励，他们使用了奖励和惩罚的语言，但他们研究的系统是适合模式识别和感知学习的监督学习系统。即使今天，一些研究者和教科书仍会弱化或模糊这些学习类型之间的区别。例如，一些人工神经网络教科书使用“试错”一词来描述从训练样例中学习的网络。这种混淆可以理解，因为这些网络使用误差信息来更新连接权重；但它错过了试错学习的本质特征，即基于评价性反馈来选择动作，而这种反馈并不依赖于知道正确动作应当是什么。

部分由于这些混淆，真正的试错学习研究在 20 世纪 60 年代和 70 年代变得少见，尽管也有一些值得注意的例外。在 20 世纪 60 年代，工程文献中首次使用“强化”和“强化学习”这两个术语来描述试错学习的工程用途（例如 Waltz and Fu, 1965；Mendel, 1966；Fu, 1970；Mendel and McClaren, 1970）。特别有影响的是 Minsky 的论文“迈向人工智能的步骤”（Minsky, 1961）；该文讨论了若干与试错学习有关的问题，包括预测、期望，以及他称为复杂强化学习系统的基本信用分配问题：怎样把成功的功劳分配给可能参与产生成功的许多决策？本书讨论的所有方法，从某种意义上说，都是朝着解决这一问题的方向努力。Minsky 的这篇论文今天仍很值得阅读。

接下来几段，我们讨论在 20 世纪 60 年代和 70 年代，对真正试错学习的计算性和理论性研究相对被忽视时，一些其他例外和部分例外。

一个例外是新西兰研究者 John Andreae 的工作。他发展了一个名为 STeLLA 的系统，它通过与环境相互作用进行试错学习。该系统包括一个世界内部模型，后来还包括一个“内部独白”，以处理隐藏状态问题（Andreae, 1963, 1969a,b）。Andreae 后来的工作（1977）更强调向教师学习，但仍然包括试错学习，并把生成新奇事件作为系统目标之一。这项工作的一个特征是“回漏过程”，Andreae（1998）更充分地阐述了它；该过程实现了一种信用分配机制，类似于我们所描述的向后更新操作。遗憾的是，他的开创性研究并不广为人知，也没有对后来的强化学习研究产生很大影响。最近已有综述可参考（Andreae, 2017a,b）。

更有影响的是 Donald Michie 的工作。1961 年和 1963 年，他描述了一个简单的试错学习系统，用于学习如何下井字棋（或圈叉棋），称为 MENACE（Matchbox Educable Naughts and Crosses Engine，可教育火柴盒圈叉引擎）。它为每个可能的游戏局面准备一个火柴盒，每个火柴盒中含有若干彩色珠子；从该局面出发，每种可能走法对应一种不同颜色。通过从与当前游戏局面对应的火柴盒中随机抽取一颗珠子，就可以确定 MENACE 的走法。一盘棋结束后，会根据 MENACE 决策的好坏，对下棋过程中使用过的盒子添加或移除珠子，以奖励或惩罚这些决策。Michie 和 Chambers（1968）描述了另一个名为 GLEE（Game Learning Expectimaxing Engine，游戏学习期望极大化引擎）的井字棋强化学习器，以及一个名为 BOXES 的强化学习控制器。他们把 BOXES 用于学习如何根据失败信号来平衡铰接在可移动小车上的杆；这个失败信号只在杆倒下或小车到达轨道末端时出现。这个任务改编自 Widrow 和 Smith（1964）较早的工作；后者使用监督学习方法，并假设由一个已经能平衡杆的教师给出指导。Michie 和 Chambers 的杆平衡版本，是在不完整知识条件下强化学习任务的最佳早期例子之一。它影响了后来的许多强化学习工作，从我们自己的若干研究开始（Barto, Sutton, and Anderson, 1983；Sutton, 1984）。Michie 始终强调试错和学习作为人工智能基本方面的作用（Michie, 1974）。

Widrow、Gupta 和 Maitra（1973）修改了 Widrow 和 Hoff（1960）的最小均方（LMS）算法，产生了一条可以从成功和失败信号而不是训练样例中学习的强化学习规则。他们称这种学习形式为“选择性自举适应”，并把它描述为“与评论家学习”而不是“与教师学习”。他们分析了这条规则，并展示了它如何学会玩二十一点。这是 Widrow 对强化学习的一次孤立探索；他的监督学习贡献更有影响力。我们使用的术语“评论家”就来自 Widrow、Gupta 和 Maitra 的论文。Buchanan、Mitchell、Smith 和 Johnson（1978）在机器学习语境中独立使用了 critic 这个术语（另见 Dietterich and Buchanan, 1984），但对他们而言，评论家是一个能够做出超出评价表现之外工作的专家系统。

学习自动机研究对通向现代强化学习研究的试错线索产生了更直接的影响。这些方法用于求解一种非联结的、纯选择性的学习问题；类比于老虎机或“单臂赌博机”，该问题称为 $k$ 臂赌博机，只不过有 $k$ 个拉杆（见第 2 章）。学习自动机是简单、低记忆的机器，用来提高这些问题中获得奖励的概率。学习自动机起源于 20 世纪 60 年代俄罗斯数学家和物理学家 M. L. Tsetlin 及其同事的工作（Tsetlin, 1973 在其身后出版），此后在工程领域得到广泛发展（见 Narendra and Thathachar, 1974, 1989）。这些发展包括随机学习自动机的研究；它们是根据奖励信号更新动作概率的方法。尽管 Harth 和 Tzanakou（1974）的 Alopex 算法（用于模式提取的算法）不是在随机学习自动机传统中发展出来的，但它是一种检测动作与强化之间相关性的随机方法，并影响了我们早期的一些研究（Barto, Sutton, and Brouwer, 1981）。随机学习自动机也被心理学中更早的工作预示过，从 William Estes（1950）为学习的统计理论所做的努力开始，并由其他人进一步发展（例如 Bush and Mosteller, 1955；Sternberg, 1963）。

心理学中发展的统计学习理论被经济学研究者采用，导致该领域出现一条致力于强化学习的研究线索。这项工作始于 1973 年，当时 Bush 和 Mosteller 的学习理论被应用于一组经典经济模型（Cross, 1973）。这一研究的一个目标，是研究行为更像真实人类而非传统理想化经济主体的人工主体（Arthur, 1991）。这种方法扩展为在博弈论语境中研究强化学习。经济学中的强化学习很大程度上独立于人工智能中早期的强化学习工作而发展，尽管博弈论仍然是两个领域都感兴趣的主题（超出本书范围）。Camerer（2011）讨论了经济学中的强化学习传统；Nowe、Vrancx 和 De Hauwere（2012）从本书介绍方法的多智能体扩展角度，概述了这一主题。博弈论语境中的强化，与用于下井字棋、跳棋和其他娱乐性游戏程序中的强化学习，是相当不同的主题。关于强化学习和游戏的这一方面，可参见 Szita（2012）的综述。

John Holland（1975）概述了一种基于选择原则的自适应系统一般理论。他的早期工作主要关心非联结形式的试错，例如演化方法和 $k$ 臂赌博机。1976 年以及 1986 年更完整地，他引入了分类器系统；这是真正的强化学习系统，包括联结和价值函数。Holland 分类器系统中的一个关键组成部分，是用于信用分配的“桶队算法”，它与我们在井字棋例子中使用、并将在第 6 章讨论的时序差分算法密切相关。另一个关键组成部分是遗传算法，这是一种演化方法，其作用是演化出有用的表示。分类器系统已由许多研究者广泛发展，形成强化学习研究的一个主要分支（Urbanowicz and Moore, 2009 综述）；但遗传算法本身，我们并不认为它单独就是强化学习系统，它受到的关注要多得多，其他演化计算方法也是如此（例如 Fogel, Owens and Walsh, 1966；Koza, 1992）。

在人工智能内部，最应归功于复兴强化学习试错线索的人是 Harry Klopf（1972, 1975, 1982）。Klopf 认识到，随着学习研究者几乎完全转向监督学习，自适应行为的一些基本方面正在丢失。按照 Klopf 的说法，缺失的是行为的享乐方面，即从环境中实现某种结果的驱力，朝着期望目标控制环境并远离不期望目标的驱力（见第 15.9 节）。这是试错学习的本质思想。Klopf 的思想对本书作者特别有影响，因为我们对这些思想的评价（Barto and Sutton, 1981a）使我们认识到监督学习与强化学习之间的区别，并最终把注意力集中到强化学习上。我们和同事完成的许多早期工作，目的都是表明强化学习和监督学习确实不同（Barto, Sutton, and Brouwer, 1981；Barto and Sutton, 1981b；Barto and Anandan, 1985）。其他研究展示了强化学习如何处理人工神经网络学习中的重要问题，特别是如何为多层网络产生学习算法（Barto, Anderson, and Sutton, 1982；Barto and Anderson, 1985；Barto, 1985, 1986；Barto and Jordan, 1987；见第 15.10 节）。

现在我们转向强化学习历史中的第三条线索，即涉及时序差分学习的线索。时序差分学习方法的独特之处在于，它们由同一量在时间上相继估计之间的差来驱动，例如井字棋例子中的获胜概率。这条线索比另外两条更小，也更不清晰，但它在该领域中发挥了特别重要的作用，部分原因是时序差分方法似乎是强化学习中新出现的独特方法。

时序差分学习的起源部分在动物学习心理学中，特别是在次级强化物这一概念中。次级强化物是与食物或疼痛等初级强化物配对过的刺激，因此具有了类似的强化性质。Minsky（1954）可能是第一个意识到这一心理学原则可能对人工学习系统重要的人。Arthur Samuel（1959）是第一个提出并实现一种包含时序差分思想的学习方法的人；这是他著名的跳棋程序的一部分（第 16.2 节）。

Samuel 没有提到 Minsky 的工作，也没有提到与动物学习的可能联系。他的灵感显然来自 Claude Shannon（1950）的建议，即可以给计算机编程，让它使用评价函数下棋，并且它也许能够通过在线修改这一函数来提高棋艺。（Shannon 的这些思想也可能影响了 Bellman，但我们没有这方面证据。）Minsky（1961）在他的“Steps”论文中广泛讨论了 Samuel 的工作，并指出其与自然和人工的次级强化理论之间的联系。

正如我们已经讨论过的，在 Minsky 和 Samuel 工作之后的十年里，几乎没有试错学习方面的计算工作，显然也完全没有时序差分学习方面的计算工作。1972 年，Klopf 把试错学习与时序差分学习的一个重要组成部分结合起来。Klopf 感兴趣的是能够扩展到大型系统学习的原则，因此他被局部强化的概念所吸引；根据这种概念，整体学习系统的子部件可以彼此强化。他发展了“广义强化”的思想：每个组成部分（名义上每个神经元）都用强化术语来看待其所有输入，把兴奋性输入看作奖励，把抑制性输入看作惩罚。这与我们现在所知的时序差分学习并不是同一个思想；回顾起来，它比 Samuel 的工作离时序差分学习更远。另一方面，Klopf 把这一思想与试错学习联系起来，并把它关联到动物学习心理学的大量经验数据库。

Sutton（1978a,b,c）进一步发展了 Klopf 的思想，特别是它们与动物学习理论的联系，描述了由时间上相继预测的变化所驱动的学习规则。他和 Barto 细化了这些思想，并基于时序差分学习发展出一个经典条件反射的心理学模型（Sutton and Barto, 1981a；Barto and Sutton, 1982）。随后出现了其他几个有影响力的、基于时序差分学习的经典条件反射心理学模型（例如 Klopf, 1988；Moore et al., 1986；Sutton and Barto, 1987, 1990）。当时发展的一些神经科学模型，也很适合用时序差分学习来解释（Hawkins and Kandel, 1984；Byrne, Gingrich, and Baxter, 1990；Gelperin, Hopfield, and Tank, 1985；Tesauro, 1986；Friston et al., 1994），尽管在大多数情况下并没有历史联系。

我们早期关于时序差分学习的工作，强烈受到动物学习理论和 Klopf 工作的影响。与 Minsky 的“Steps”论文以及 Samuel 跳棋程序之间的关系，只是在后来才被认识到。不过，到 1981 年，我们已经充分意识到，上面提到的所有早期工作都是时序差分和试错两条线索的一部分。此时，我们发展出一种把时序差分学习与试错学习结合起来的方法，称为 actor--critic 架构，并把它应用于 Michie 和 Chambers 的杆平衡问题（Barto, Sutton, and Anderson, 1983）。这一方法在 Sutton（1984）的博士论文中得到广泛研究，并在 Anderson（1986）的博士论文中扩展为使用反向传播神经网络。大约在这一时期，Holland（1986）以桶队算法的形式，明确把时序差分思想并入他的分类器系统。Sutton（1988）迈出了关键一步：把时序差分学习从控制中分离出来，把它作为一种一般预测方法来处理。那篇论文还引入了 TD($\lambda$) 算法，并证明了它的一些收敛性质。

1981 年，当我们正在完成关于 actor--critic 架构的工作时，我们发现 Ian Witten（1977, 1976a）的一篇论文；它似乎是时序差分学习规则最早的出版物。他提出的方法就是我们现在称为表格型 TD(0) 的方法，用作求解 MDP 的自适应控制器的一部分。这项工作最初于 1974 年提交给期刊出版，也出现在 Witten 1976 年的博士论文中。Witten 的工作是 Andreae 早期 STeLLA 实验和其他试错学习系统的后代。因此，Witten 1977 年的论文横跨强化学习研究的两条主要线索，即试错学习和最优控制，同时也对时序差分学习做出了独特的早期贡献。

时序差分和最优控制这两条线索在 1989 年因 Chris Watkins 发展出 Q-learning 而完全结合起来。这项工作扩展并整合了强化学习研究三条线索中的先前工作。Paul Werbos（1987）自 1977 年以来一直主张试错学习和动态规划的融合，并对这一整合作出了贡献。到 Watkins 工作出现时，强化学习研究已经有了巨大增长，主要发生在人工智能的机器学习子领域，但也发生在人工神经网络以及更广义的人工智能中。1992 年，Gerry Tesauro 的双陆棋程序 TD-Gammon 获得显著成功，为该领域带来了更多关注。

自本书第一版出版以来，神经科学中发展出了一个繁荣的子领域，关注强化学习算法与神经系统中强化学习之间的关系。对此负主要责任的，是时序差分算法的行为与大脑中产生多巴胺的神经元活动之间一种惊人的相似性；这一点由许多研究者指出（Friston et al., 1994；Barto, 1995a；Houk, Adams, and Barto, 1995；Montague, Dayan, and Sejnowski, 1996；Schultz, Dayan, and Montague, 1997）。第 15 章将介绍强化学习的这一令人兴奋的方面。强化学习近代历史中的其他重要贡献太多，无法在这份简短叙述中一一提及；我们会在它们出现的各章末尾引用其中许多。

\section*{书目说明}

关于强化学习的更多一般性介绍，我们建议读者参考 Szepesvari（2010）、Bertsekas and Tsitsiklis（1996）、Kaelbling（1993a）以及 Sugiyama、Hachiya 和 Morimura（2013）的著作。采取控制或运筹学视角的书包括 Si、Barto、Powell 和 Wunsch（2004）、Powell（2011）、Lewis 和 Liu（2012）以及 Bertsekas（2012）。Cao（2009）的综述把强化学习放在随机动态系统的其他学习与优化方法的语境中。期刊 \textit{Machine Learning} 的三期专刊聚焦强化学习：Sutton（1992a）、Kaelbling（1996）和 Singh（2002）。有用的综述由 Barto（1995b）、Kaelbling、Littman 和 Moore（1996）以及 Keerthi 和 Ravindran（1997）提供。Wiering 和 van Otterlo（2012）编辑的文集对近期发展提供了出色概述。

\begin{longtable}{p{0.12\textwidth}p{0.78\textwidth}}
1.2 & 本章中 Phil 早餐的例子受到 Agre（1988）的启发。 \\
1.5 & 井字棋例子中使用的时序差分方法将在第 6 章中展开。 \\
\end{longtable}

\end{document}
